\documentclass[../../main.tex]{subfiles}
\begin{document}

\section{光滑流形的例子}\label{section:ch01-光滑流形的例子}

在进一步讨论一般理论之前，让我们浏览一些光滑流形的例子。

\begin{example}[$0$ 维流形]\label{example:ch01-0维流形}
维数为 $0$ 的拓扑流形 $M$ 恰好就是一个可数离散空间。对每个点 $p\in M$，$p$ 的唯一同胚于 $\mathbb{R}^{0}$ 中开子集的邻域就是 $\{p\}$ 本身，而恰好存在一个坐标映射 $\varphi:\{p\}\to\mathbb{R}^{0}$。因此，$M$ 上所有图所成集合平凡地满足光滑相容条件，且每个 $0$ 维流形都有唯一的光滑结构。
\end{example}

\begin{example}[欧几里得空间]\label{example:ch01-欧几里得空间}
对每个非负整数 $n$，欧几里得空间 $\mathbb{R}^{n}$ 是一个光滑 $n$ 流形，其光滑结构由单个图 $(\mathbb{R}^{n},\operatorname{Id}_{\mathbb{R}^{n}})$ 组成的图册所确定。我们称这个为 $\mathbb{R}^{n}$ 上的标准光滑结构，称所得坐标映射为标准坐标。除非我们另行明确说明，我们总在 $\mathbb{R}^{n}$ 上使用这个光滑结构。关于这个光滑结构，$\mathbb{R}^{n}$ 的光滑坐标图恰好是那些使 $\varphi$ 从 $U$ 到另一个开子集 $\widetilde{U}\subset\mathbb{R}^{n}$ 为（通常微积分意义下）微分同胚的图 $(U,\varphi)$。
\end{example}

\begin{example}[$\mathbb{R}$ 上的另一个光滑结构]\label{example:ch01-另一光滑结构}
考虑由
\begin{align}
\psi(x)=x^{3}\label{align:ch01-另一光滑结构}
\end{align}
给出的同胚 $\psi:\mathbb{R}\to\mathbb{R}$。由单个图 $(\mathbb{R},\psi)$ 组成的图册在 $\mathbb{R}$ 上定义了一个光滑结构。这个图与标准光滑结构不是光滑相容的，因为转移映射 $\operatorname{Id}_{\mathbb{R}}\circ\psi^{-1}(y)=y^{1/3}$ 在原点不光滑。因此，由 $\psi$ 在 $\mathbb{R}$ 上定义的光滑结构与标准的不同。利用类似的想法，不难在任意给定的正维数拓扑流形上构造许多不同的光滑结构，只要它一开始有一个光滑结构即可（见习题~1-6）。
\end{example}

\begin{example}[有限维向量空间]\label{example:ch01-有限维向量空间}
设 $V$ 是一个有限维实向量空间。$V$ 上的任意范数都确定一个拓扑，该拓扑与范数的选取无关（练习~B.49）。带这个拓扑，$V$ 是一个拓扑 $n$ 流形，并具有如下定义的自然光滑结构。$V$ 的每个（有序）基 $(E_{1},\dots,E_{n})$ 定义一个基同构 $E:\mathbb{R}^{n}\to V$，由
\begin{align*}
E(x)=\sum_{i=1}^{n}x^{i}E_{i}
\end{align*}
给出。这个映射是同胚，所以 $(V,E^{-1})$ 是一个图。若 $(\widetilde{E}_{1},\dots,\widetilde{E}_{n})$ 是任意另一个基，而 $\widetilde{E}(x)=\sum_{j}x^{j}\widetilde{E}_{j}$ 是相应的同构，则存在某个可逆矩阵 $(A_{i}^{j})$ 使对每个 $i$ 有 $E_{i}=\sum_{j}A_{i}^{j}\widetilde{E}_{j}$。于是这两个图之间的转移映射由 $\widetilde{E}^{-1}\circ E(x)=\widetilde{x}$ 给出，其中 $\widetilde{x}=(\widetilde{x}^{1},\dots,\widetilde{x}^{n})$ 由
\begin{align*}
\sum_{j=1}^{n}\widetilde{x}^{j}\widetilde{E}_{j}=\sum_{i=1}^{n}x^{i}E_{i}=\sum_{i,j=1}^{n}x^{i}A_{i}^{j}\widetilde{E}_{j}
\end{align*}
确定。由此可得 $\widetilde{x}^{j}=\sum_{i}A_{i}^{j}x^{i}$。因此，把 $x$ 映到 $\widetilde{x}$ 的映射是一个可逆线性映射，因而是微分同胚，所以任意两个这样的图都光滑相容。所有这类图所成集合由此定义了一个光滑结构，称为 $V$ 上的标准光滑结构。
\end{example}

\subsubsection{Einstein 求和约定}\label{subsubsection:ch01-Einstein求和约定}
这里正好暂停一下，介绍一个在光滑流形研究中常用的重要记号约定。由于本主题中像 $\sum_{i}x^{i}E_{i}$ 这样的求和大量出现，我们常通过略去求和号来缩写这种和，如
\begin{align*}
E(x)=x^{i}E_{i},
\end{align*}
作为 $E(x)=\sum_{i=1}^{n}x^{i}E_{i}$ 的缩写。我们按如下规则解释任何这样的表达式，该规则称为 Einstein 求和约定：如果同一个指标名（如上面表达式中的 $i$）在任意单项项中恰好出现两次，一次作为上指标、一次作为下指标，则该项被理解为对该指标的所有可能取值求和，一般从 $1$ 到所讨论空间的维数。这个简单想法由 Einstein 引入，旨在通过消除显式写出求和号的必要来降低光滑流形研究中出现的表达式的复杂度。我们在全书中系统地使用求和约定（附录除外，许多读者会先读附录再读正文）。

求和约定的另一个重要方面是指标的位置。我们总把基向量（如 $E_{i}$）写成带下指标，把向量关于一组基的分量（如 $x^{i}$）写成带上指标。这些指标约定有助于保证：在具有数学意义的求和中，每个需求和的指标在任意给定项中通常出现两次，一次作为下指标、一次作为上指标。任何被隐含求和的指标都是"哑指标"，意思是把每个哑指标替换为不同的名字，该表达式的值不变。例如，$x^{i}E_{i}$ 与 $x^{j}E_{j}$ 表示完全相同的东西。

由于点 $(x^{1},\dots,x^{n})\in\mathbb{R}^{n}$ 的坐标也就是它关于标准基的分量，为与我们把向量分量写成带上指标的约定一致，我们需要对这些坐标使用上指标，全书都这样做。虽然乍看这可能别扭，但与求和约定结合后，在处理复杂的带指标和时它提供了巨大优势，其中之一是：不具数学意义的表达式往往会因违反指标约定而很快暴露自己。（主要例外是涉及欧几里得点积 $x\cdot y=\sum_{i}x^{i}y^{i}$ 的表达式，其中同一指标两次出现在上方位置；以及 $\mathbb{R}^{2n}$ 上的标准辛形式，我们将在第~22 章定义。在这类表达式中我们总是显式写出求和号。）

\subsubsection{更多例子}\label{subsubsection:ch01-更多例子}
现在我们继续给出光滑流形的例子。

\begin{example}[矩阵空间]\label{example:ch01-矩阵空间}
令 $M(m\times n;\mathbb{R})$ 表示元素为实数的 $m\times n$ 矩阵所成集合。因为它在矩阵加法和数乘下是一个维数为 $mn$ 的实向量空间，$M(m\times n;\mathbb{R})$ 是一个光滑的 $mn$ 维流形。（事实上，常常有用的是把 $M(m\times n;\mathbb{R})$ 等同于 $\mathbb{R}^{mn}$，只需把所有矩阵元排成一行。）类似地，$m\times n$ 复矩阵空间 $M(m\times n;\mathbb{C})$ 是 $\mathbb{R}$ 上维数为 $2mn$ 的向量空间，因而是一个维数为 $2mn$ 的光滑流形。在 $m=n$（方阵）的特殊情形，我们把 $M(n\times n;\mathbb{R})$ 和 $M(n\times n;\mathbb{C})$ 分别缩写为 $M(n;\mathbb{R})$ 和 $M(n;\mathbb{C})$。
\end{example}

\begin{example}[开子流形]\label{example:ch01-开子流形}
设 $U$ 是 $\mathbb{R}^{n}$ 的任意开子集。则 $U$ 是一个拓扑 $n$ 流形，而单个图 $(U,\operatorname{Id}_{U})$ 在 $U$ 上定义了一个光滑结构。

更一般地，设 $M$ 是一个光滑 $n$ 流形，$U\subset M$ 是任意开子集。由
\begin{align*}
\mathcal{A}_{U}=\{\text{$M$ 的光滑图 $(V,\varphi)$ 满足 }V\subset U\}
\end{align*}
定义 $U$ 上的一个图册。每个点 $p\in U$ 都包含在 $M$ 的某个图 $(W,\varphi)$ 的定义域中；若令 $V=W\cap U$，则 $(V,\varphi|_{V})$ 是 $\mathcal{A}_{U}$ 中一个定义域包含 $p$ 的图。因此，$U$ 被 $\mathcal{A}_{U}$ 中图的定义域覆盖，且容易验证这是 $U$ 的一个光滑图册。所以 $M$ 的任意开子集本身自然地是一个光滑 $n$ 流形。赋予这个光滑结构后，我们把任意开子集称为 $M$ 的一个开子流形。（我们将在第~5 章定义一类更一般的子流形。）
\end{example}

\begin{example}[一般线性群]\label{example:ch01-一般线性群}
一般线性群 $GL(n;\mathbb{R})$ 是元素为实数的可逆 $n\times n$ 矩阵所成集合。它是一个光滑的 $n^{2}$ 维流形，因为它是 $n^{2}$ 维向量空间 $M(n;\mathbb{R})$ 的一个开子集，即（连续的）行列式函数非零的那个集合。
\end{example}

\begin{example}[满秩矩阵]\label{example:ch01-满秩矩阵}
前面的例子对满秩长方矩阵有一个自然推广。设 $m<n$，令 $M_{m}(m\times n;\mathbb{R})$ 表示 $M(m\times n;\mathbb{R})$ 中由秩为 $m$ 的矩阵组成的子集。若 $A$ 是任意这样的矩阵，则 $\operatorname{rank}A=m$ 意味着 $A$ 有某个非奇异的 $m\times m$ 子矩阵。由行列式函数的连续性，同一个子矩阵在 $M(m\times n;\mathbb{R})$ 中 $A$ 的一个邻域上行列式非零，这意味着 $A$ 有一个包含于 $M_{m}(m\times n;\mathbb{R})$ 的邻域。因此，$M_{m}(m\times n;\mathbb{R})$ 是 $M(m\times n;\mathbb{R})$ 的一个开子集，因而本身是一个光滑的 $mn$ 维流形。类似的论证表明，当 $n<m$ 时 $M_{n}(m\times n;\mathbb{R})$ 是一个光滑的 $mn$ 维流形。
\end{example}

\begin{example}[线性映射空间]\label{example:ch01-线性映射空间}
设 $V$ 和 $W$ 是有限维实向量空间，令 $L(V;W)$ 表示从 $V$ 到 $W$ 的线性映射所成集合。因为 $L(V;W)$ 本身是一个有限维向量空间（其维数是 $V$ 与 $W$ 维数之积），它具有如例~\cref{example:ch01-有限维向量空间} 中的自然光滑流形结构。在它上面放置整体坐标的一种方式是为 $V$ 和 $W$ 选取基，并把每个 $T\in L(V;W)$ 用其矩阵表示，这给出 $L(V;W)$ 与 $M(m\times n;\mathbb{R})$ 的一个同构，其中 $m=\dim W$，$n=\dim V$。
\end{example}

\begin{example}[光滑函数的图像]\label{example:ch01-光滑函数图像}
若 $U\subset\mathbb{R}^{n}$ 是一个开子集，$f:U\to\mathbb{R}^{k}$ 是一个光滑函数，我们在前面（例~\cref{example:ch01-连续函数的图像}）已经看到，$f$ 的图像在子空间拓扑下是一个拓扑 $n$ 流形。由于 $\Gamma(f)$ 被单个图像坐标图 $\varphi:\Gamma(f)\to U$（$\pi_{1}$ 的限制）覆盖，我们可以通过宣告图像坐标图 $(\Gamma(f),\varphi)$ 为光滑图，在 $\Gamma(f)$ 上放置一个典范光滑结构。
\end{example}

\begin{example}[球面]\label{example:ch01-球面光滑结构}
我们在例~\cref{example:ch01-球面} 中证明了 $S^{n}\subset\mathbb{R}^{n+1}$ 是一个拓扑 $n$ 流形。我们按如下方式在 $S^{n}$ 上放置光滑结构。对每个 $i=1,\dots,n+1$，令 $(U_{i}^{\pm},\varphi_{i}^{\pm})$ 表示我们在例~\cref{example:ch01-球面} 中构造的图像坐标图。对任意不同的指标 $i$ 和 $j$，转移映射 $\varphi_{i}^{\pm}\circ(\varphi_{j}^{\pm})^{-1}$ 容易计算。在 $i<j$ 的情形，我们得到
\begin{align*}
\varphi_{i}^{\pm}\circ(\varphi_{j}^{\pm})^{-1}(u^{1},\dots,u^{n})
=(u^{1},\dots,\widehat{u^{i}},\dots,\pm\sqrt{1-|u|^{2}},\dots,u^{n})
\end{align*}
（其中平方根在第 $j$ 个位置），当 $i>j$ 时有类似公式。当 $i=j$ 时，一个更简单的计算给出
\begin{align*}
\varphi_{i}^{+}\circ(\varphi_{i}^{-})^{-1}=\varphi_{i}^{-}\circ(\varphi_{i}^{+})^{-1}=\operatorname{Id}_{B^{n}}.
\end{align*}
因此，图族 $\{(U_{i}^{\pm},\varphi_{i}^{\pm})\}$ 是一个光滑图册，从而在 $S^{n}$ 上定义了一个光滑结构。我们称其为 $S^{n}$ 的标准光滑结构。
\end{example}

\begin{example}[水平集]\label{example:ch01-水平集}
前面的例子可推广如下。设 $U\subset\mathbb{R}^{n}$ 是一个开子集，$\Phi:U\to\mathbb{R}$ 是一个光滑函数。对任意 $c\in\mathbb{R}$，集合 $\Phi^{-1}(c)$ 称为 $\Phi$ 的一个水平集。选某个 $c\in\mathbb{R}$，令 $M=\Phi^{-1}(c)$，并假设恰好对每个 $a\in\Phi^{-1}(c)$ 全导数 $D\Phi(a)$ 非零。因为 $D\Phi(a)$ 是一个行矩阵，其元素为偏导数 $(\partial\Phi/\partial x^{1}(a),\dots,\partial\Phi/\partial x^{n}(a))$，对每个 $a\in M$ 存在某个 $i$ 使 $\partial\Phi/\partial x^{i}(a)\neq 0$。由隐函数定理（定理~C.40，其中 $x^{i}$ 扮演 $y$ 的角色）可得，存在 $a$ 的一个邻域 $U_{0}$，使 $M\cap U_{0}$ 可以表示为形如
\begin{align*}
x^{i}=f\bigl(x^{1},\dots,\widehat{x^{i}},\dots,x^{n}\bigr)
\end{align*}
的方程的图像，其中 $f$ 是定义在 $\mathbb{R}^{n-1}$ 的某个开子集上的光滑实值函数。因此，正如 $n$ 球面的情形那样论证，我们看到 $M$ 是一个维数为 $(n-1)$ 的拓扑流形，并具有一个光滑结构，使上面这种选取 $f$ 所得到的每个图像坐标图都是光滑图。在第~5 章中，我们将发展光滑子流形理论，它是这一构造的深远推广。
\end{example}

\begin{example}[射影空间]\label{example:ch01-射影空间光滑结构}
$n$ 维实射影空间 $\mathbb{RP}^{n}$ 由例~\cref{example:ch01-射影空间} 是一个拓扑 $n$ 流形。让我们验证该例中构造的坐标图 $(U_{i},\varphi_{i})$ 全都光滑相容。为方便起见假设 $i>j$，直接计算得
\begin{align*}
\varphi_{j}\circ\varphi_{i}^{-1}(u^{1},\dots,u^{n})
=\left(\frac{u^{1}}{u^{j}},\dots,\frac{u^{j-1}}{u^{j}},\frac{u^{j+1}}{u^{j}},\dots,\frac{u^{i-1}}{u^{j}},\frac{1}{u^{j}},\frac{u^{i}}{u^{j}},\dots,\frac{u^{n}}{u^{j}}\right),
\end{align*}
这是从 $\varphi_{i}(U_{i}\cap U_{j})$ 到 $\varphi_{j}(U_{i}\cap U_{j})$ 的微分同胚。
\end{example}

\begin{example}[光滑乘积流形]\label{example:ch01-光滑乘积流形}
若 $M_{1},\dots,M_{k}$ 分别是维数为 $n_{1},\dots,n_{k}$ 的光滑流形，我们在例~\cref{example:ch01-乘积流形} 中证明了乘积空间 $M_{1}\times\cdots\times M_{k}$ 是一个维数为 $n_{1}+\cdots+n_{k}$ 的拓扑流形，其坐标图形如 $(U_{1}\times\cdots\times U_{k},\varphi_{1}\times\cdots\times\varphi_{k})$。任意两个这样的图都光滑相容，因为容易验证
\begin{align*}
(\psi_{1}\times\cdots\times\psi_{k})\circ(\varphi_{1}\times\cdots\times\varphi_{k})^{-1}
=(\psi_{1}\circ\varphi_{1}^{-1})\times\cdots\times(\psi_{k}\circ\varphi_{k}^{-1}),
\end{align*}
这是一个光滑映射。这在乘积上定义了一个自然光滑流形结构，称为乘积光滑流形结构。例如，这在 $n$ 环面 $T^{n}=S^{1}\times\cdots\times S^{1}$ 上给出一个光滑流形结构。
\end{example}

在我们迄今所见的每个例子中，都分两个阶段构造光滑流形结构：从一个拓扑空间出发，验证它是拓扑流形，然后指定一个光滑结构。把这两步合并为一个单一构造通常更方便，尤其是当我们从一个尚未配备拓扑的集合出发时。下面的引理提供了一条捷径——它表明，给定一个带有若干"光滑重叠"的适当"坐标图"的集合，我们如何能用这些图在该集合上同时定义拓扑和光滑结构。

\begin{lemma}[光滑流形坐标图引理]\label{lemma:ch01-光滑流形坐标图引理}
设 $M$ 是一个集合，并给定 $M$ 的一个子集族 $\{U_{\alpha}\}$ 以及映射 $\varphi_{\alpha}:U_{\alpha}\to\mathbb{R}^{n}$，满足以下性质：
\begin{enumerate}[(i)]
\item 对每个 $\alpha$，$\varphi_{\alpha}$ 是 $U_{\alpha}$ 与开子集 $\varphi_{\alpha}(U_{\alpha})\subset\mathbb{R}^{n}$ 之间的双射。
\item 对每个 $\alpha$ 和 $\beta$，集合 $\varphi_{\alpha}(U_{\alpha}\cap U_{\beta})$ 与 $\varphi_{\beta}(U_{\alpha}\cap U_{\beta})$ 在 $\mathbb{R}^{n}$ 中是开的。
\item 每当 $U_{\alpha}\cap U_{\beta}\neq\varnothing$ 时，映射 $\varphi_{\beta}\circ\varphi_{\alpha}^{-1}:\varphi_{\alpha}(U_{\alpha}\cap U_{\beta})\to\varphi_{\beta}(U_{\alpha}\cap U_{\beta})$ 光滑。
\item 可数多个集合 $U_{\alpha}$ 覆盖 $M$。
\item 每当 $p,q$ 是 $M$ 中两个不同点时，或者存在某个 $U_{\alpha}$ 同时包含 $p$ 和 $q$，或者存在不相交的集合 $U_{\alpha},U_{\beta}$ 使 $p\in U_{\alpha}$ 且 $q\in U_{\beta}$。
\end{enumerate}
则 $M$ 具有唯一的光滑流形结构，使每个 $(U_{\alpha},\varphi_{\alpha})$ 都是光滑图。
\end{lemma}

\begin{figure}[ht]
\centering
\includegraphics[width=0.8\textwidth]{fig-01-9.png}
\caption{光滑流形图引理}\label{figure:ch01-1.9}
\end{figure}

\begin{proof}
我们通过取所有形如 $\varphi_{\alpha}^{-1}(V)$（其中 $V$ 是 $\mathbb{R}^{n}$ 的开子集）的集合为基来定义拓扑。为证明这是一个拓扑的基，我们需要证明：对于两个基集合 $\varphi_{\alpha}^{-1}(V)$ 和 $\varphi_{\beta}^{-1}(W)$ 的交中的任意点 $p$，存在第三个基集合包含 $p$ 且包含于该交中。只需证明 $\varphi_{\alpha}^{-1}(V)\cap\varphi_{\beta}^{-1}(W)$ 本身是一个基集合即可（图~\cref{figure:ch01-1.9}）。为此注意 (iii) 蕴涵 $(\varphi_{\beta}\circ\varphi_{\alpha}^{-1})^{-1}(W)$ 是 $\varphi_{\alpha}(U_{\alpha}\cap U_{\beta})$ 的一个开子集，而 (ii) 蕴涵该集合在 $\mathbb{R}^{n}$ 中也是开的。由此可得
\begin{align*}
\varphi_{\alpha}^{-1}(V)\cap\varphi_{\beta}^{-1}(W)
=\varphi_{\alpha}^{-1}\bigl(V\cap(\varphi_{\beta}\circ\varphi_{\alpha}^{-1})^{-1}(W)\bigr)
\end{align*}
也是一个基集合，如所断言。

每个映射 $\varphi_{\alpha}$ 于是是到其像的同胚（本质上由定义），所以 $M$ 是 $n$ 维局部欧几里得的。Hausdorff 性质由 (v) 容易得到，第二可数性由 (iv) 及练习~A.22 的结果得到，因为每个 $U_{\alpha}$ 是第二可数的。最后，(iii) 保证族 $\{(U_{\alpha},\varphi_{\alpha})\}$ 是一个光滑图册。显然，这个拓扑和光滑结构是满足引理结论的唯一一对。
\end{proof}

\begin{example}[Grassmann 流形]\label{example:ch01-Grassmann流形}
设 $V$ 是一个 $n$ 维实向量空间。对任意整数 $0\leqslant k\leqslant n$，令 $G_{k}(V)$ 表示 $V$ 的所有 $k$ 维线性子空间所成的集合。我们将证明 $G_{k}(V)$ 可以自然地赋予一个维数为 $k(n-k)$ 的光滑流形结构。赋予这个结构后，它称为 Grassmann 流形，或简称 Grassmann 流形（Grassmannian）。在 $V=\mathbb{R}^{n}$ 的特殊情形，Grassmann 流形 $G_{k}(\mathbb{R}^{n})$ 常记作某些更简单的记号，如 $G_{k,n}$ 或 $G(k;n)$。注意 $G_{1}(\mathbb{R}^{n+1})$ 恰好就是 $n$ 维射影空间 $\mathbb{RP}^{n}$。

在 $G_{k}(V)$ 上构造光滑结构比我们迄今所做的那些稍许复杂些，但基本想法不过是用线性代数为 $G_{k}(V)$ 构造坐标图，然后应用光滑流形坐标图引理。我们将在第~21 章给出 $G_{k}(V)$ 是光滑流形的一个更简短证明（见例~21.21）。

设 $P$ 和 $Q$ 是 $V$ 的任意一对互补子空间，维数分别为 $k$ 和 $n-k$，于是 $V$ 分解为直和 $V=P\oplus Q$。任意线性映射 $X:P\to Q$ 的图像可以等同于一个 $k$ 维子空间 $\Gamma(X)\subset V$，定义为
\begin{align*}
\Gamma(X)=\{v+Xv:v\in P\}.
\end{align*}
任意这样的子空间都具有与 $Q$ 的交为零子空间的性质。反之，任意与 $Q$ 平凡相交的子空间 $S\subset V$ 都是某个唯一线性映射 $X:P\to Q$ 的图像，它可构造如下：令 $\pi_{P}:V\to P$ 和 $\pi_{Q}:V\to Q$ 为直和分解所确定的投影；于是假设蕴涵 $\pi_{P}|_{S}$ 是从 $S$ 到 $P$ 的同构。因此，$X=(\pi_{Q}|_{S})\circ(\pi_{P}|_{S})^{-1}$ 是一个良定义的从 $P$ 到 $Q$ 的线性映射，且容易验证 $S$ 是它的图像。

令 $L(P;Q)$ 表示从 $P$ 到 $Q$ 的线性映射所成向量空间，并令 $U_{Q}$ 表示 $G_{k}(V)$ 中与 $Q$ 的交为平凡的那些 $k$ 维子空间所成的子集。对应 $X\mapsto\Gamma(X)$ 定义一个映射 $\Gamma:L(P;Q)\to U_{Q}$，上面的讨论表明 $\Gamma$ 是一个双射。令 $\varphi=\Gamma^{-1}:U_{Q}\to L(P;Q)$。通过为 $P$ 和 $Q$ 选取基，我们可以把 $L(P;Q)$ 等同于 $M((n-k)\times k;\mathbb{R})$，进而等同于 $\mathbb{R}^{k(n-k)}$，于是可以把 $(U_{Q},\varphi)$ 看作一个坐标图。由于每个这样的图的像都是整个 $L(P;Q)$，引理~\cref{lemma:ch01-光滑流形坐标图引理} 的条件 (i) 显然满足。

现在取任意另一对子空间 $(P',Q')$，令 $\pi_{P'},\pi_{Q'}$ 为相应的投影，$\varphi':U_{Q'}\to L(P';Q')$ 为相应映射。集合 $\varphi(U_{Q}\cap U_{Q'})\subset L(P;Q)$ 由所有图像与 $Q'$ 平凡相交的线性映射 $X:P\to Q$ 组成。为看出这个集合在 $L(P;Q)$ 中是开的，对每个 $X\in L(P;Q)$，令 $I_{X}:P\to V$ 为映射 $I_{X}(v)=v+Xv$，它是从 $P$ 到 $X$ 的图像的双射。因为 $\Gamma(X)=\operatorname{Im}I_{X}$ 且 $Q'=\operatorname{Ker}\pi_{P'}$，由练习~B.22(d) 可得，$X$ 的图像与 $Q'$ 平凡相交当且仅当 $\pi_{P'}\circ I_{X}$ 满秩。因为 $\pi_{P'}\circ I_{X}$ 的矩阵元（关于任意基）连续依赖于 $X$，例~\cref{example:ch01-满秩矩阵} 的结果表明所有这类 $X$ 所成集合在 $L(P;Q)$ 中是开的。于是光滑流形坐标图引理中的性质 (ii) 成立。

我们需要证明转移映射 $\varphi'\circ\varphi^{-1}$ 在 $\varphi(U_{Q}\cap U_{Q'})$ 上光滑。任取 $X\in\varphi(U_{Q}\cap U_{Q'})\subset L(P;Q)$，令 $S$ 表示子空间 $\Gamma(X)\subset V$。若令 $X'=\varphi'\circ\varphi^{-1}(X)$，则如前所述 $X'=(\pi_{Q'}|_{S})\circ(\pi_{P'}|_{S})^{-1}$（图~\cref{figure:ch01-1.10}）。为把这个映射与 $X$ 联系起来，注意 $I_{X}:P\to S$ 是同构，所以我们可以写
\begin{align*}
X'=(\pi_{Q'}|_{S})\circ I_{X}\circ(I_{X})^{-1}\circ(\pi_{P'}|_{S})^{-1}
=(\pi_{Q'}\circ I_{X})\circ(\pi_{P'}\circ I_{X})^{-1}.
\end{align*}

\begin{figure}[ht]
\centering
\includegraphics[width=0.8\textwidth]{fig-01-10.png}
\caption{$G_{k}(V)$ 上坐标的光滑相容性}\label{figure:ch01-1.10}
\end{figure}

为证明这光滑依赖于 $X$，定义线性映射 $A:P\to P'$、$B:P\to Q'$、$C:Q\to P'$、$D:Q\to Q'$ 如下：
\begin{align*}
A=\pi_{P'}|_{P},\qquad B=\pi_{Q'}|_{P},\qquad C=\pi_{P'}|_{Q},\qquad D=\pi_{Q'}|_{Q}.
\end{align*}
于是对 $v\in P$，有
\begin{align*}
(\pi_{P'}\circ I_{X})v=(A+CX)v,\qquad (\pi_{Q'}\circ I_{X})v=(B+DX)v,
\end{align*}
由此可得 $X'=(B+DX)(A+CX)^{-1}$。一旦我们为 $P,Q,P',Q'$ 选取基，所有这些线性映射都由矩阵表示。因为由 Cramer 法则，$(A+CX)^{-1}$ 的矩阵元是 $A+CX$ 的矩阵元的有理函数，由此可得 $X'$ 的矩阵元光滑依赖于 $X$ 的矩阵元。这证明了 $\varphi'\circ\varphi^{-1}$ 是光滑映射，所以我们构造的图满足引理~\cref{lemma:ch01-光滑流形坐标图引理} 的条件 (iii)。

为检验条件 (iv)，我们只需注意 $G_{k}(V)$ 事实上可以被有限多个集合 $U_{Q}$ 覆盖：例如，若 $(E_{1},\dots,E_{n})$ 是 $V$ 的任意固定基，把基元素分成含 $k$ 个和 $n-k$ 个元素的两个子集的任意划分都确定适当的子空间 $P$ 和 $Q$，而任意子空间 $S$ 对至少其中一个划分都与 $Q$ 平凡相交（见练习~B.9）。因此，$G_{k}(V)$ 被由固定基的所有可能划分所确定的有限多个图覆盖。

最后，Hausdorff 条件 (v) 容易验证，只需注意对任意两个 $k$ 维子空间 $P,P'\subset V$，可以找到一个维数为 $n-k$ 的子空间 $Q$，使它与 $P$ 和 $P'$ 的交都平凡，于是 $P$ 和 $P'$ 都包含在由（比方说）$(P,Q)$ 所确定的图的定义域中。
\end{example}

\end{document}